\section*{第 \ref{ch:b2:frontier-orbitals} 章——前线轨道与反应活性}

\begin{solution}{exo:b2:frontier-orbitals:1}
乙烯：HOMO 为 $\alpha + \beta$，LUMO 为 $\alpha - \beta$。烯丙基负离子（四个电子）：HOMO 为 $\alpha$（非键
能级），LUMO 为 $\alpha - 1.414\beta$。丁二烯：HOMO 为 $\alpha + 0.618\beta$，LUMO 为 $\alpha -
0.618\beta$。
\end{solution}

\begin{solution}{exo:b2:frontier-orbitals:2}
硬试剂：\ce{OH-}、\ce{H+}、\ce{F-}。软试剂：\ce{I-}、\ce{RS-}、碘甲烷的碳原子。
\end{solution}

\begin{solution}{exo:b2:frontier-orbitals:3}
丁二烯（可取 s-顺式构象）和环戊二烯（锁定为 s-顺式）能够反应；
1,4-戊二烯不共轭，不能反应。(2\textit{Z},4\textit{Z})-二烯要达到
s-顺式构象，只能使两个内侧甲基相互挤压，因此它
反应性很差。
\end{solution}

\begin{solution}{exo:b2:frontier-orbitals:4}
环己-3-烯-1-甲腈：六元环，双键位于二烯体原来的 C2
和 C3 之间，\ce{C#N} 基团位于由
亲双烯体形成的两个环碳之一上。
\end{solution}

\begin{solution}{exo:b2:frontier-orbitals:5}
$\Delta\varepsilon = \qty{8}{eV}$：微扰法得 $2\beta^2/\Delta\varepsilon = \qty{0.25}{eV}$。精确解中较低能级为
$-6 - \sqrt{16 + 1} = \qty{-10.123}{eV}$，能量降低 $2 \times 0.123 = \qty{0.246}{eV}$。
\end{solution}

\begin{solution}{exo:b2:frontier-orbitals:6}
碘甲烷是软亲电试剂：受轨道控制，由具有较大
HOMO 系数的原子（碳）进攻。自由碳正离子是硬亲电试剂：受电荷控制，
部分由氮原子进攻，氮原子带有更多负电荷。
\end{solution}

\begin{solution}{exo:b2:frontier-orbitals:7}
二烯体的 C4（HOMO 系数最大）与丙烯醛的 $\beta$ 碳成键，该碳的
LUMO 系数最大：生成 2-甲氧基环己-3-烯-1-甲醛，两个取代基位于
相邻碳原子上（“邻位”产物）。
\end{solution}

\begin{solution}{exo:b2:frontier-orbitals:8}
带有酯基的双环[2.2.1]庚烯（降冰片烯）骨架，酯基连在
原亲双烯体的一个碳上，朝向新生成的 \ce{C=C} 键，位于六元环之下
（内型），在与 \ce{CH2} 桥相反的一侧。
\end{solution}

\begin{solution}{exo:b2:frontier-orbitals:9}
两个羰基降低了亲双烯体的 LUMO：它与环戊二烯 HOMO 之间的能隙
远小于乙烯的情形，过渡态的稳定化能 $\beta^2/\Delta\varepsilon$ 更大，
因此反应能垒更低。
\end{solution}

\begin{solution}{exo:b2:frontier-orbitals:10}
该反应是立体专一的：马来酸酯生成两个酯基均为顺式的加合物
（双内型为主要产物；它是非手性的，即内消旋化合物）；富马酸酯生成
反式加合物，一个酯基为内型、一个为外型，以一对对映异构体（外消旋）的形式生成。
\end{solution}

\begin{solution}{exo:b2:frontier-orbitals:11}
两个 $\pi$ 键变为两个更强的 $\sigma$ 键：$\Delta_r H^\circ < 0$。两个分子变为一个：
$\Delta_r S^\circ < 0$。$\Delta_r G^\circ = \Delta_r H^\circ - T\Delta_r S^\circ$ 随 $T$ 增大，并在
$T_i = \Delta_r H^\circ/\Delta_r S^\circ$ 以上变为正值：高温下加合物会重新分解
（逆第尔斯–阿尔德反应）。
\end{solution}

\begin{solution}{exo:b2:frontier-orbitals:12}
两个充满电子的孤电子对轨道并排时形成四电子相互作用，这种相互作用
是排斥的。分子绕 \ce{N-N} 键转动，直到两个孤电子对
大致垂直（为邻位交叉构象），此时它们的重叠与排斥作用
都很小。
\end{solution}

\begin{solution}{pb:b2:frontier-orbitals:1}
\textbf{1.} 五元环，含两个共轭双键和一个 \ce{CH2}：该环
使二烯单元保持 s-顺式构象。
\textbf{2.} 亲双烯体，通过其中一个双键参与反应。
\textbf{3.} 降冰片烯骨架与一个环戊烯环稠合；两个新的 $\sigma$ 键将
二烯体的两端与亲双烯体双键上的两个碳相连。
\textbf{4.} 内型加合物（内型规则）：亲双烯体上的其余环在过渡态中位于
二烯体的下方。
\textbf{5.} 是的：该反应为协同反应，两个新键均在各反应物的同一面形成，并且
亲双烯体的构型保持不变。
\textbf{6.} 环戊二烯既是良好的二烯体（锁定为 s-顺式），也是亲双烯体，并且
该反应为协同反应，没有需要稳定的带电中间体。
\textbf{7.} 负值：两个 $\pi$ 键被两个 $\sigma$ 键取代。
\textbf{8.} 负值：两个分子生成一个。
\textbf{9.} $\Delta_r G^\circ = \Delta_r H^\circ - T\Delta_r S^\circ$，两项均为负：当 $T$ 较低时为负，
在 $T_i = \Delta_r H^\circ/\Delta_r S^\circ$ 时为零，高于此温度为正。
\textbf{10.} 室温下 $T < T_i$：二聚化有利，但速度慢。在
二聚体沸点附近，平衡有利于单体，
且很快达到。
\textbf{11.} 单体以蒸气形式离开热的混合物：移去生成物使
平衡持续向单体方向移动（勒夏特列）。
\textbf{12.} 室温下它再次二聚；低温会减慢该反应。
\textbf{13.} $0.62 - (-0.62) = 1.24|\beta|$.
\textbf{14.} $0.62 - (-0.35) = 0.97|\beta|$.
\textbf{15.} 环戊二烯与马来酸酐：较小的能隙使过渡态
的稳定化程度更大，反应更快。
\textbf{16.} 它们与 \ce{C=C} 键共轭且具有吸电子性：与
丙烯醛一样，它们将 $\pi^*$ 能级拉低。
\textbf{17.} 其 LUMO 在羰基碳上具有波瓣，这些羰基碳朝向
二烯体 HOMO 的中心碳。
\textbf{18.} 降冰片烯骨架；酸酐环位于新双键下方，与
\ce{CH2} 桥相对。
\textbf{19.} 两个桥头碳与带有酸酐基团的两个碳：共
四个手性碳原子。
\textbf{20.} 否：一个镜面穿过 \ce{CH2} 桥，并位于分子
两半之间；它是内消旋化合物。
\textbf{21.} 它们在亲双烯体上为顺式，协同加成保持它们为顺式。
\textbf{22.} 骨架相同，但酸酐环位于 \ce{CH2} 桥一侧。
\textbf{23.} 否：将内型改为外型，需要断裂并重新形成 \ce{C-C} 键，
这须经由需要加热的逆反应完成。
\textbf{24.} 水使酸酐开环：生成顺式二元羧酸，两个 \ce{COOH} 基团均为内型。
\textbf{25.} 内型加合物有 $\boldsymbol{4}$ 个手性碳原子。
\end{solution}
